% ID: SEL-D0021-OA-E02
% ESTADO: cotejada. Segunda pasada visual completa y borrador compilado revisados.
% Fuente: Selectividad/matematicas_aplicadas/reserva_a_EXAMEN 3-ANDALUCÍA.pdf
% SHA-256: 0b99df3e2c9a7af0a074136b72c3d257ad9241ccc79fed32894de424481b6be2
% Página: 1; opción A; ejercicio original 2.

El beneficio en euros que obtiene una empresa al vender $x$ unidades de un artículo viene dado por la función $B(x)=-x^2+360x-18000$, $50\leq x\leq 350$.
\begin{enumerate}[label=\alph*)]
\item \textbf{(0.8 puntos)} \textquestiondown{}Cuál es el beneficio obtenido si vende 100 unidades? \textquestiondown{}Cuántas unidades debe vender para obtener un beneficio de 13500 \texteuro{}?
\item \textbf{(1 punto)} \textquestiondown{}Cuál es el número de unidades que debe vender para que el beneficio sea máximo? \textquestiondown{}A cuánto asciende ese beneficio?
\item \textbf{(0.7 puntos)} Represente gráficamente la función y determine cuántas unidades hay que vender para no obtener pérdidas.
\end{enumerate}
